\bookchapter{Глава 63. Представления, состояние и навигация}

\chapterintro{Делать игрушки из одного экрана вы уже умеете. Эта глава решает две проблемы «повзросления»: как перемещаться между несколькими экранами (навигация) и как данные текут между экранами (управление состоянием). Пройдя её, вы сможете собрать классический скелет многоэкранного приложения «список детали редактирование».}

\booksection{«Панорама потоков данных» в управлении состоянием}

У SwiftUI не один инструмент состояния — не только @State. Сначала взгляните на панораму и разберитесь, «когда что использовать»:

\begin{center}
\begin{tabularx}{\linewidth}{|X|X|X|X|}
\hline
\textbf{Инструмент} & \textbf{За что отвечает} & \textbf{Кто использует} & \textbf{Типичный сценарий} \tabularnewline
\hline
@State & временное состояние, приватное для экрана & текущее представление & счётчик, текст в поле ввода \tabularnewline
@Binding & заимствование чужого состояния & дочернее представление меняет родительское & свои компоненты \tabularnewline
@StateObject / @Observable & логика, общая для нескольких экранов & модель представления (ViewModel) & загрузка данных, бизнес-логика \tabularnewline
@EnvironmentObject / .environment & глобальное общее состояние & всё приложение & статус входа пользователя, настройки темы \tabularnewline
\hline
\end{tabularx}
\end{center}

\begin{center}\small\color{gray}Панорама инструментов состояния SwiftUI (2026 год, @Observable уже мейнстрим)\end{center}

\begin{lstlisting}[language=Swift]
// Современный способ на iOS 17+: макрос @Observable (рекомендуется для новых проектов)
import Observation

@Observable
class CounterModel {                 // модель представления
    var count = 0
    func increment() { count += 1 }
}

struct CounterView: View {
    @State private var model = CounterModel()   // хранит модель

    var body: some View {
        VStack {
            Text("Вы нажали \(model.count) раз")
            Button("Нажми") { model.increment() }
        }
    }
}
\end{lstlisting}

Грубый совет новичку: для приватного состояния экрана — @State, для межэкранного/с логикой — классы @Observable. @StateObject/@ObservableObject из старых учебников — вчерашний стиль: достаточно уметь читать, а новый код пишите, как сказано выше.

\booksection{NavigationStack: «стек» переходов между экранами}

Начиная с iOS 16, навигация по нескольким экранам строится на NavigationStack — как стопка тарелок: кладёшь одну сверху, при возврате снимаешь по одной:

\begin{lstlisting}[language=Swift]
struct HomeView: View {
    var body: some View {
        NavigationStack {
            List {
                NavigationLink("Записная книжка Ивана") {
                    DetailView(name: "Записная книжка Ивана")   // целевой экран
                }
                NavigationLink("Записная книжка Марии") {
                    DetailView(name: "Записная книжка Марии")
                }
            }
            .navigationTitle("Список записных книжек")   // крупный заголовок сверху
        }
    }
}

struct DetailView: View {
    let name: String
    var body: some View {
        Text("Подробности: \(name)")
            .navigationTitle(name)
            .navigationBarTitleDisplayMode(.inline)
    }
}
\end{lstlisting}

NavigationLink — «кликабельная ссылка-переход»; navigationTitle автоматически создаёт заголовок панели навигации и кнопку «назад». Когда нужен программный переход (например, автоматический возврат после нажатия «Сохранить»), закрывайте текущий экран через @Environment(\textbackslash{}.dismiss) — этих двух приёмов хватает на 90\% навигационных задач.

\booksection{Всплывающие окна и нижние вкладки: sheet и TabView}

Кроме «проталкивания экранов», есть ещё два классических приёма взаимодействия:

\begin{lstlisting}[language=Swift]
// 1. Выезжает снизу на пол/весь экран (любимый приём для форм ввода)
struct HomeView: View {
    @State private var showAdd = false
    var body: some View {
        Button("Новая запись") { showAdd = true }
            .sheet(isPresented: $showAdd) {
                AddRecordView()     // появляющаяся форма ввода
            }
    }
}

// 2. Нижняя панель вкладок ("скелет" приложения)
struct MainTabView: View {
    var body: some View {
        TabView {
            RecordListView()
                .tabItem { Label("Записи", systemImage: "book") }
            StatsView()
                .tabItem { Label("Статистика", systemImage: "chart.pie") }
            SettingsView()
                .tabItem { Label("Настройки", systemImage: "gear") }
        }
    }
}
\end{lstlisting}

\begin{tcbinfo}{Совет: когда что использовать?}
Эмпирическое правило: основной поток строим на NavigationStack (список детали), временные действия — во всплывающем окне sheet (добавление, правка, подтверждение), несколько равноправных крупных разделов — в TabView (вкладки внизу). Вместе эти три элемента — скелет подавляющего большинства приложений.

\end{tcbinfo}

\booksection{List: «народный список» iOS}

Самый ходовой и самый «яблочный» компонент в iOS — List: разделители, удаление свайпом и подсветка нажатия идут у него из коробки:

\begin{lstlisting}[language=Swift]
struct RecordListView: View {
    @State private var records = ["Обед 25.5", "Метро 6.0", "Зарплата +8000"]

    var body: some View {
        List {
            ForEach(records, id: \.self) { r in
                Text(r)
            }
            .onDelete { indexSet in          // свайп влево — удалить
                records.remove(atOffsets: indexSet)
            }
        }
    }
}
\end{lstlisting}

Список + NavigationLink + sheet + TabView — соединив эти четыре элемента, вы уже способны собрать полный скелет приложения «список записей открыть детали нажать + и добавить запись». Именно так построен проект книги учёта в следующей главе.

\begin{tcbpractice}{Практика: соберите «каркас из трёх страниц»}
\booksection{Жизненный цикл приложения: передний план и фон}

Подобно жизненному циклу Activity в Android, приложение iOS проходит состояния «рождение, передний план, фон, смерть». В SwiftUI за ними наблюдают через значение окружения scenePhase:

\begin{lstlisting}[language=Swift]
@main
struct MyApp: App {
    @Environment(\.scenePhase) private var phase

    var body: some Scene {
        WindowGroup {
            ContentView()
        }
        .onChange(of: phase) { _, newPhase in
            switch newPhase {
            case .active:    print("Снова на переднем плане, возобновляем обновление")
            case .inactive:  print("Переход (звонок / мгновение в фоне)")
            case .background: print("Ушёл в фон, сохраняем черновик")
            @unknown default: break
            }
        }
    }
}
\end{lstlisting}

Типичное применение: обновить данные при возвращении на передний план, сохранить черновик при уходе в фон. iOS очень строго следит за фоновой работой — обычное приложение вскоре после ухода в фон приостанавливается («замораживается»), так что не рассчитывайте тайком делать в фоне тяжёлую работу; если нужны задачи по расписанию, используйте локальные уведомления из главы 68, а не «вечное проживание в фоне».

Соберите скелет книги учёта на фиктивных данных: TabView с тремя вкладками (Книга / Статистика / Настройки) страница книги через List показывает фиктивные данные, нажатие на строку ведёт к деталям (NavigationStack) кнопка «+» в правом верхнем углу открывает sheet для ввода. Данные не сохраняются, сеть не используется — чистая тренировка каркаса; когда скелет готов, наполнить его настоящими данными в главе 65 будет уже делом техники.

\end{tcbpractice}

\begin{tcbsummary}{Краткие итоги главы}
\begin{itemize}
  \item Четыре инструмента состояния: @State (приватно для страницы), @Binding (заимствование), @Observable (логика), .environment (глобально)
  \item В новых проектах пишите модель представления макросом @Observable; @ObservableObject из старых уроков достаточно уметь читать
  \item NavigationStack открывает страницы; @Environment(\textbackslash{}.dismiss) — программный возврат назад
  \item Основной поток — страницы, временные действия — sheet, равноправные крупные разделы — TabView
  \item List — «народный список»: разделители и удаление свайпом доступны из коробки
\end{itemize}

\end{tcbsummary}
